% =============================================================
% Bilgisayar Mimarisi — Sürüm Ayarları
% Bu dosya sürüm numarası, çalışma modu ve bölüm yönetimini
% merkezi olarak kontrol eder.
% =============================================================

% --- Sürüm Bilgisi ---
\newcommand{\kitapsurum}{0.5.0}

% --- Çalışma Modu ---
% Günlük çalışma: sadece etkin bölümü derlemek için aşağıdaki satırı aç.
% Tam derleme veya sürüm hazırlığı için yorum satırı yap.
% Günlük çalışma: sadece etkin bölümü derle (uncomment et)
%\includeonly{ekB-sayisal-tasarim/ekB}

% Yayın sürümü: sadece yayındaki bölümleri derle (uncomment et)
\includeonly{bolum01-giris/bolum01,bolum02-basarim/bolum02,bolum03-buyruk-kumesi/bolum03,bolum04-programlama/bolum04-programlama,bolum05-aritmetik/bolum05,ekA-sayi-sistemleri/ekA,ekB-sayisal-tasarim/ekB,ekC-risc-v-referans/ekC,ekD-pratik-rehber/ekD}

% --- Bölüm Yönetim Makroları ---
% \bolumekle: Dosya varsa \include eder, yoksa sayacı atlar ve label tanımlar.
%   #1 = dosya yolu (uzantısız), #2 = label adı, #3 = bölüm numarası, #4 = başlık
% \ekekle: Aynı mantık, ek bölümler için (appendix label formatı).
\makeatletter
\newcommand{\bolumekle}[4]{%
  \IfFileExists{#1.tex}{%
    \include{#1}%
  }{%
    \stepcounter{chapter}%
    \newlabel{#2}{{#3}{}{#4}{chapter.#3}{}}%
  }%
}
\newcommand{\ekekle}[4]{%
  \IfFileExists{#1.tex}{%
    \include{#1}%
  }{%
    \stepcounter{chapter}%
    \newlabel{#2}{{#3}{}{#4}{appendix.#3}{}}%
  }%
}
% Ek B içindeki section-level referanslar (dosya yoksa tanımla)
\AtBeginDocument{%
  \@ifundefined{r@sec:ekB-yazmaclar}{%
    \newlabel{sec:ekB-yazmaclar}{{B.4}{}{Yazmaçlar}{section.B.4}{}}%
  }{}%
}
\makeatother
